%% =====================================================================
%%  T-04-02  The Favour Debt
%%  -------------------------------------------------------------------
%%  One idea per file. Core is mandatory. Slots in canonical order.
%%  Max 700 words. No layout commands. No filler. New source material
%%  for this entry goes in sources/<BOOK>/T-04-02.tex -- never by
%%  rewriting this file (docs/RULES.md R0.1).
%% =====================================================================
%% @kind: topic
%% @id: T-04-02
%% @title: The Favour Debt
%% @subtitle: Gratitude is the cheapest lever and the least reliable payment
%% @chapter: c04
%% @order: 20
%% @status: stable
%% @sources: B1, B3
%% @updated: 2026-09-19

\topic{T-04-02}{The Favour Debt}{Gratitude is the cheapest lever and the least reliable payment}

\begin{core}
A favour creates an account, and both sides know it. The person who gave can draw on it later; the person who received carries an obligation that has no stated amount and no stated due date, which makes it unusually easy to over-draw.
\end{core}
\begin{mechanism}
The obligation works because refusing to repay feels like a statement about your character rather than about the transaction. That is what makes it leverage: the cost of refusing is paid in self-image, not in money. It is also why the lever is unreliable --- the account has no agreed balance, so the debtor can always conclude that the debt has been settled, or was never incurred, or was discharged by the resentment the giving produced. Both sources therefore treat gratitude as a poor instrument in either direction: weak to collect on, and expensive to owe.
\end{mechanism}
\begin{tells}
  \item A gift or favour arrives before anything has been asked, and is larger than the relationship warrants.
  \item The giver mentions what they did for you, in front of others.
  \item You find yourself unable to say no to a small request because of a large kindness.
  \item Resentment appears where gratitude was expected --- the debtor's way of cancelling the account.
\end{tells}
\begin{moves}
  \item Ask in terms of the other party's interest, not in terms of what you have done for them; reminding someone of a past favour gives them a reason to ignore you.
  \item Keep accounts small and settled quickly, so that no balance accumulates.
  \item Where you must create an obligation, make the repayment specific and immediate, which is worth more later than a vague large debt.
  \item If you have done someone a favour you intend to draw on, do not mention it. Mention converts an asset into a grievance.
\end{moves}
\begin{counters}
  \item Price the favour at the moment it arrives. A specific, prompt repayment closes the account and removes the leverage entirely.
  \item Decline unsolicited largeness. An unrequested gift is the classic opening move because it cannot be refused without rudeness --- which is exactly why refusing it politely and early is worth the awkwardness.
  \item Separate the two questions: \emph{was I helped?} and \emph{does this request follow from it?} Most attempts to draw on the account collapse when the second is asked on its own.
  \item Watch your own resentment. It is evidence that you believe you owe something you did not agree to owe.
\end{counters}
\begin{limits}
The debt is real in relationships that are meant to be ongoing --- family, long colleagues --- and treating every kindness as an opening move eventually leaves you with no one willing to do you one. Both sources accept that generosity functions as an instrument while also noting that it is one of the few instruments whose value depends on not being visibly used. The practical rule is asymmetric: be alert to accounts being opened, be relaxed about ones being closed.
\end{limits}

\sources{B1, B3}
\seesources
\seealso{T-04-01, T-18-01, T-25-02}
